\chapter*{Abstract}
\addcontentsline{toc}{chapter}{Abstract}

This thesis documents the design, implementation, and evaluation of \TimeLens{},
an AI- and AR-powered mobile guide for the \GEM{} (GEM). Pointing a phone at an
exhibit, a visitor sees the artifact recognized in real time and can ask follow-up
questions answered in English or Arabic. The work targets three problems specific
to in-gallery deployment: fine-grained visual similarity among \textbf{51 catalogued
artifacts} (many near-identical Ramesside statues), the gap between curated training
data and handheld phone-camera conditions, and the risk of an AI guide stating
historical ``facts'' it cannot support.

The thesis presents two engineering-research contributions and a market study. The
first is an \textbf{on-device artifact detector} developed through a
data-quality-driven iteration study: from foundation-model auto-annotation
(YOLO-World), through spatial label-cleaning rules, to a fully hand-annotated
dataset. The study isolates \emph{label quality}---rather than model architecture or
input resolution---as the decisive factor: the final \textbf{YOLOv8n} model resolves
every previously failing class while remaining a \textbf{5.97\,MB} TensorFlow Lite
asset that runs in real time on a mid-range phone. The second is a \textbf{bilingual
Retrieval-Augmented Generation (RAG) guide} grounded in a \textbf{108-record ChromaDB}
knowledge base; a systematic benchmark of \textbf{seven candidate language models}
selects \textbf{Gemma 4 E2B (Q4\_K\_M)}, and a ten-step optimization reduces
end-to-end latency from \textbf{over 30\,s to about 10\,s}. Both ship in a production
\textbf{Flutter} application with a bilingual interface, museum location gating, and
text-to-speech.

Beyond headline metrics, the thesis is written to be reproducible and honest about
its limits: the detection study explicitly discloses validation \emph{frame-leakage}
and the training-to-deployment \emph{domain shift}, reporting a conservative held-out
result (mAP@0.5 $=0.751$) alongside the in-distribution figures. The document
includes reproducibility notes, an ethics and bias statement, and a formal
generative-AI disclosure.

\begin{insightbox}
A reliable museum guide is built less by adding model complexity than by
\emph{grounding} it: clean labels make the on-device detector trustworthy, and
retrieval-grounded generation keeps the language model factual. \TimeLens{}'s value
comes from disciplined data and grounding, not a larger model.
\end{insightbox}
